To derive Hamilton's equations from the given Hamiltonian, we start by considering the Hamiltonian function:

$$
H(x, p) = p \dot{x} - L(x, \dot{x}).
$$

Here, \( H \) is a function of the independent variables \( x \) and \( p \), while \( L \) is a function of \( x \) and \( \dot{x} \). The relationship between \( \dot{x} \) and \( p \) is given by the Legendre transformation, where

$$
p = \frac{\partial L}{\partial \dot{x}}.
$$

To find Hamilton's equations, we need to compute the partial derivatives of \( H \) with respect to \( x \) and \( p \).

**1. Derivative of \( H \) with respect to \( x \):**

Using the chain rule, we have:

$$
\frac{\partial H}{\partial x} = \frac{\partial}{\partial x} (p \dot{x} - L) = -\frac{\partial L}{\partial x}.
$$

Since the Euler-Lagrange equation gives:

$$
\frac{d}{dt} \left( \frac{\partial L}{\partial \dot{x}} \right) = \frac{\partial L}{\partial x},
$$

we substitute \( \frac{\partial L}{\partial \dot{x}} = p \) to get:

$$
\frac{d}{dt}(p) = \frac{\partial L}{\partial x}.
$$

Thus,

$$
\frac{\partial H}{\partial x} = -\dot{p}.
$$

**2. Derivative of \( H \) with respect to \( p \):**

For the partial derivative with respect to \( p \), we have:

$$
\frac{\partial H}{\partial p} = \frac{\partial}{\partial p} (p \dot{x} - L) = \dot{x},
$$

since \( \dot{x} \) is treated as a constant with respect to \( p \).

Therefore, Hamilton's equations are:

$$
\frac{\partial H}{\partial x} = -\dot{p}, \quad \frac{\partial H}{\partial p} = \dot{x}.
$$